\documentclass[12pt,letterpaper,oneside]{article}
\input{ITESCA/maestria-en-gestion-administrativa/plan-de-negocios-mga/actividades-generadas/2026-II/revision-2026-09-20/formato-itesca-plan-negocios.tex}
\begin{document}
\portadaITESCA{Preguntas para el análisis de demanda y oferta}{6548}
Este documento es un borrador hipotético pendiente de confirmación. Presenta preguntas e ítems propuestos para analizar la demanda y la oferta de AM Taller Autocentro, proyecto local planteado como taller, llantera y refaccionaria. Los instrumentos no han sido aplicados; por tanto, no se incluyen respuestas, resultados ni conclusiones de mercado. La zona de estudio, la población, la muestra y el periodo de referencia permanecen pendientes de definición.

\section{Propósito de la recolección}

El propósito es recopilar información que permita fundamentar decisiones sobre ventas, capacidad de prestación de servicios, promoción, ubicación y precio. El análisis se divide en dos poblaciones: usuarios de vehículos que podrían requerir servicios automotrices y establecimientos o proveedores relacionados con mecánica, llantas y refacciones.

Para estudiar la demanda se propone una encuesta dirigida a personas que poseen, administran o utilizan vehículos. Para la oferta se plantea una entrevista semiestructurada con responsables de establecimientos y proveedores potenciales, o una lista de cotejo limitada a información pública, observable y autorizada. No se presupone que deban aplicarse simultáneamente los tres tipos de instrumentos.

\section{Listado para el análisis de la demanda}

\subsection{Uso del vehículo y necesidades de atención}

Las siguientes preguntas constituyen una propuesta de encuesta. Antes de aplicarla deberán establecerse opciones de respuesta claras, un periodo de referencia y un aviso de participación voluntaria. También deberá evitarse la recopilación de datos personales que no sean indispensables.

\begin{enumerate}
    \item ¿Con qué frecuencia utiliza un vehículo para actividades personales, familiares o laborales?
    \item ¿Qué tipo general de vehículo utiliza con mayor frecuencia y cuál es su antigüedad aproximada?
    \item ¿Quién toma habitualmente la decisión de contratar servicios o comprar refacciones para ese vehículo?
    \item ¿Con qué frecuencia solicita mantenimiento preventivo, diagnóstico o reparación mecánica?
    \item Durante el periodo definido para el estudio, ¿qué servicios automotrices ha requerido?
    \item ¿En qué circunstancias suele buscar atención: mantenimiento programado, falla inesperada, revisión preventiva, cambio de llantas u otra necesidad?
    \item ¿Qué dificultades ha encontrado al solicitar servicios de mecánica, llantas o refacciones?
\end{enumerate}

\subsection{Elección del establecimiento}

\begin{enumerate}
    \item ¿Qué factores considera más importantes al elegir un taller, llantera o refaccionaria?
    \item ¿Qué importancia asigna a la confianza, ubicación, tiempo de atención, precio, garantía y disponibilidad?
    \item ¿Qué medios utiliza para localizar y comparar establecimientos automotrices?
    \item ¿Qué información necesita conocer antes de autorizar un servicio o comprar una refacción?
    \item ¿Prefiere recibir una cotización desglosada antes de autorizar el trabajo? ¿Qué conceptos espera que incluya?
    \item ¿Qué canal de comunicación prefiere para solicitar información y recibir seguimiento?
    \item ¿Consideraría útil contar con una agenda para seleccionar una fecha de revisión? Explique su respuesta.
    \item ¿Qué días y horarios le resultarían más convenientes para acudir a un establecimiento?
    \item ¿Qué tiempo de traslado consideraría aceptable para recibir el servicio?
    \item ¿Qué elementos aumentarían su confianza al solicitar una primera cita?
\end{enumerate}

\subsection{Precio, promoción e intención de compra}

\begin{enumerate}
    \item ¿Cómo determina si el precio de un servicio automotriz es razonable?
    \item Al comparar alternativas, ¿qué importancia tiene el precio frente a la garantía, la claridad de la cotización y el tiempo estimado de atención?
    \item ¿Qué información promocional le ayudaría a conocer un nuevo establecimiento automotriz?
    \item ¿Por qué medio preferiría recibir información sobre servicios y recomendaciones de mantenimiento?
    \item ¿Estaría dispuesto a considerar los servicios propuestos por AM Taller Autocentro si cumplieran sus criterios de atención, ubicación, precio y claridad? ¿Por qué?
\end{enumerate}

Las respuestas permitirían identificar necesidades, preferencias y criterios de compra, pero no demostrarían por sí solas una demanda efectiva. Si se incorporan rangos de gasto o precio, deberán definirse después de una revisión preliminar y una prueba del cuestionario, sin presentar tarifas inexistentes ni inducir respuestas.

\section{Listado para el análisis de la oferta}

\subsection{Establecimientos automotrices}

Los siguientes ítems pueden utilizarse en una entrevista semiestructurada o adaptarse como criterios de una lista de cotejo. La observación deberá limitarse a información verificable, sin registrar información confidencial.

\begin{enumerate}
    \item ¿Qué categorías de servicios comunica u ofrece el establecimiento?
    \item ¿Se especializa en mecánica, llantas o refacciones, o combina varias categorías?
    \item ¿Qué horarios de atención comunica al público?
    \item ¿Atiende mediante cita, recepción directa o ambas modalidades?
    \item ¿Qué canales utiliza para proporcionar información, cotizaciones y seguimiento?
    \item ¿Qué información comunica sobre diagnóstico, autorización, tiempos estimados y garantías?
    \item ¿Cómo presenta sus precios: cotización individual, tarifa visible, paquete o valoración previa?
    \item ¿Qué atributos utiliza para diferenciar su propuesta?
    \item ¿Qué medios promocionales emplea y qué tipo de contenido difunde?
    \item ¿La ubicación presenta identificación visible, acceso y condiciones para recibir vehículos?
    \item ¿Qué servicios tienen mayor presencia en su comunicación comercial?
    \item ¿Qué limitaciones de horario, capacidad o disponibilidad declara?
\end{enumerate}

La presencia de un servicio en la publicidad no deberá interpretarse como evidencia de ventas, demanda o capacidad efectiva. Tampoco se atribuirán características que no puedan comprobarse.

\subsection{Proveedores potenciales}

\begin{enumerate}
    \item ¿Qué categorías de llantas, refacciones o insumos puede cotizar el proveedor?
    \item ¿Qué requisitos solicita para establecer una relación comercial?
    \item ¿Cuáles son sus condiciones de cotización, pedido mínimo, pago, entrega, devolución y garantía?
    \item ¿Cómo confirma la compatibilidad de una referencia antes de surtirla?
    \item ¿Con qué frecuencia actualiza su catálogo, disponibilidad y condiciones comerciales?
    \item ¿Qué documentos proporciona para identificar los productos y respaldar una operación?
    \item ¿Qué procedimiento aplica para atender errores, devoluciones o reclamaciones?
\end{enumerate}

Estos ítems no presuponen contratos, existencias, precios, compatibilidades ni acuerdos con proveedores. Las referencias de catálogo deberán considerarse opciones por cotizar hasta contar con confirmación documental.

\section{Uso de la información y control de calidad}

La información de demanda podría orientar la selección de servicios prioritarios, los horarios, los canales de contacto, los mensajes promocionales, los criterios de ubicación y la presentación de cotizaciones. El análisis de la oferta permitiría comparar modalidades de atención, capacidad operativa comunicada, categorías de servicio y condiciones potenciales de abastecimiento. Las decisiones de ventas, producción del servicio, promoción, ubicación y precio no deberán sustentarse en opiniones aisladas.

Antes de aplicar los instrumentos será necesario verificar que las preguntas sean neutrales, comprensibles y pertinentes; definir opciones de respuesta suficientes; realizar una prueba de comprensión; establecer un método de codificación; y documentar la fecha, el medio de aplicación y el criterio de selección de participantes.

Para convertir este borrador en un trabajo real se requiere confirmar el número mínimo de ítems, el formato, la rúbrica y el peso individual de la actividad; delimitar la zona geográfica; definir la población, el tamaño y método de selección de la muestra; elegir el instrumento definitivo para estudiar la oferta; establecer el periodo de referencia; y precisar las categorías de respuesta. Solo después de una aplicación autorizada podrán incorporarse resultados, interpretaciones y conclusiones.
\clearpage
\printbibliography[title={Fuentes de consulta}]
\end{document}
